\documentclass[10pt]{article}

\usepackage{parskip}
\usepackage[margin=1in]{geometry} 
\usepackage{amsmath,amsthm,amssymb, graphicx, multicol, array}
\usepackage{makecell}
\usepackage{enumitem}
\usepackage{amssymb}
\usepackage{float}
\usepackage{href-ul}
 
\newcommand{\N}{\mathbb{N}}
\newcommand{\Z}{\mathbb{Z}}
 
\newenvironment{problem}[2][Problem]{\begin{trivlist}
\item[\hskip \labelsep {\bfseries #1}\hskip \labelsep {\bfseries #2.}]}{\end{trivlist}}

\begin{document}
 
\title{Planning models cont.}
\author{Jakob Sverre Alexandersen\\
GRA4154 Supply Chain Optimization with Mathematical Programming\\
Lecture 7}
\maketitle

\tableofcontents
\newpage

\section{Problem 1}

A company produces two products sharing the same production capacity. 

The production manager plans production based on a forecast of how much is going to be shipped out during a given horizon. 

Some production capacity is lost due to setup times (changeover times) if both products are produced in the same week. 

The production manager tries to produce only one product per week, but often needs to produce both products. 

Products are stored in inventory silos. Each product has a dedicated inventory capacity

Demand must be covered, and production and inventory capacities must be respected 

Build a mathematical planning model. The objective function should be to maximize the total value of production, minus inventory costs. 

\textbf{Data}: 

Current forecast:
\begin{center}
    \begin{tabular}{|c|c|c|c|c|c|c|c|c|c|}
        \hline 
        & week 1& week 2& week 3& week 4& week 5& week 6& week 7& week 8& week 9\\\hline 
        \textbf{Product 1} & 6500 & 5900 & 6500 & 7700 & 8300 & 7700 & 5800 & 7700 & 8600 \\\hline 
        \textbf{Product 2} & 4800 & 4700 & 5900 & 8200 & 16,400 & 5000 & 13,000 & 12,500 & 11,000 \\\hline 
    \end{tabular}
\end{center}

Production capacity is 20,000 tons per week if one product is produced during the week. If both products are produced, the total capacity is reduced to 16,000 tons 
\begin{center}
    \begin{tabular}{|c|c|c|c|c|}
        \hline 
        & value per ton & inventory cost /t/w & init inventory & inventory cap. \\\hline 
        \textbf{Product 1} & 240 & 8 & 3700 & 9000\\\hline 
        \textbf{Product 2} & 285 & 9 & 6500 & 10,500\\\hline 
    \end{tabular}
\end{center}

\textbf{Sets}
\begin{itemize}
    \item $P = \{1, 2\} $: products
    \item $T = \{1, \dots, 9\} $: set of weeks 
\end{itemize}
\textbf{Parameters}
\begin{itemize}
    \item $d_{p, t} $: demand for product $p $ at time $t $
    \item $v_p $: value per ton of product $ p $
    \item $h_p $: inventory cost per ton per week for product $p $
    \item $I_p^0 $: intial inventory of product $p $
    \item $\bar{I}_p $: inventory capacity for product $p $
\end{itemize}
\textbf{Decision variables}
\begin{itemize}
    \item $x_{p, t} \ge 0$: tons of product $p $ produced in week $t $
    \item $I_{p, t} \ge 0 $: inventory of product $p $ at the end of week $t $
    \item $y_{p, t} \in \{0, 1\} $: 1 if product $p $ is produced in week $t $, 0 otherwise 
    \item $z_t \in \{0, 1\} $: 1 if both products are produced in week $t $, 0 otherwise
\end{itemize}
\newpage 
\textbf{Model}
\begin{gather*}
    \min Z = \sum_{p\in P}\sum_{t\in T}v_p x_{p, t} - \sum_{p\in P}\sum_{t\in T}h_pI_{p, t}\\
    \text{s.t.}\\
    I_{p, t} = I_{p, t - 1} + x_{p, t} - d_{p, t}\quad\forall\, p\in P,\, t\in T,\, I_{p, 0} = I_p^0\\
    I_{p, t}\le \bar{I}_p\quad\forall p\in P,\, t\in T\\
    x_{p, t} \le C^\text{single}y_{p, t}\quad\forall p\in P,\, t\in T\\
    z_t \ge y_{1, t} + y_{2, t} - 1\quad\forall t\in T\\
    \sum_{p\in P}\le 20,000 - 4,000 z_t \quad\forall t\in T\\
\end{gather*}
\section{Static planning model}
\textbf{Basic model: the economic order quantity model (EOQ)}
\begin{itemize}
    \item Assumptions: 
    \begin{itemize}
        \item Demand is constant and continuous -- $D$ per year 
        \item The lot size (decision variable) is $Q $
        \item A fixed cost (order cost or setup cost) -- $A $ -- is incurred every time a replenishment (order) takes place 
    \end{itemize}
    \item Example: 
    \begin{itemize}
        \item $D = 1000 $ per year 
        \item $A = 200 $ per replenishment
        If the lot size $Q = 100 $, the number of replenishments per year = $D / Q = 1000 / 100 = 10$

        Total order costs per year $ = 10\times 200 = 2,000 $
    \end{itemize}
    \item To minimize these fixed costs: maximize $Q $, to get as few replenishments as possible 
    
Inventory cost is assumed to be an annual inventory rate, $r $, multiplied with the value of the product, $v $

Assume $r = 0.1 \text{ and } v  = 500 $

Inventory holding cost per year for $Q = 100 $: 
\begin{gather*}
    v\times r\times Q / 2 = 500\times 0.1\times 100/2 = 2500 
\end{gather*}
(remember: 100/2 = 50 is the average inventory for $Q = 100 $)
\end{itemize}
\newpage 
\section{Finding the optimal trade-off}
between order costs and inventory holding costs 

\textbf{Total relevant costs (TRC)}: total order costs + total inventory holding costs 
\begin{gather*}
    = A\times D/Q + v\times r\times Q/2 
\end{gather*}
In the example above, we assume $Q = 100,\, r = 0.1 \text{ and } v  = 500 $, TRC = 
\begin{gather*}
    200\times 1000 / 100 + 500\times 0.2\times 100/2\\
    = 2000 + 2500 = 4500
\end{gather*}
But which level of $Q $ will minimize TCR?
\begin{gather*}
    \text{TCR}'_Q = -\frac{AD}{Q^2} + \frac{vr}{2} \to -\frac{AD}{Q^2} + \frac{vr}{2} = 0\\
    Q = \sqrt{\frac{2AD}{vr}}
\end{gather*}
(this is the ``famous'' EOQ formula)
\section{The EOQ formula with quantity discounts}
\textbf{Data}:
\begin{itemize}
    \item Demand per year $D = 1000 $
    \item Ordering cost $A = 200 $
    \item Value per unit $v = 500 $
    \item Interest rate per year $r = 0.1 $
\end{itemize}
Assume the cost price will be 499 instead of 500 if the order quantity is at least 150 units (the breakpoint). Remember, the previous EOQ was 84.44

\begin{itemize}
    \item Let $v_0 = 500 $ be the price without the quantity discount
    \item and $v_0\times (1 - d) = 499 $ be the price with quantity discount
    \item hence, $d $ is the percentage discount, in this case $(500 - 499) / 500 = 0.02 $
\end{itemize}
\textbf{The TRC when there are quantity discounts}
\begin{itemize}
    \item The purchasing price $v $ must now be included in the TRC 
    \item Without discount: TRC = $AD / Q + Qv_0r / 2+ Dv_0 $
    \item With discount: TRC = $AD / Q + Qv_0(1 - d)r / 2+ Dv_0 (1 - d)$

    (but only for $Q \ge Q_b $ where $b $ is the breakpoint)
\end{itemize}
\end{document}